\documentclass[11pt]{article}

\usepackage[T1]{fontenc}
\usepackage{amsmath,amssymb,graphicx}
\usepackage[margin=1in]{geometry}
\usepackage{booktabs}
\usepackage{hyperref}

\title{Full Combinatorial Austerity Closure of the No-Distance C Line}
\author{Tomislav Rupi\'c \\ Independent Researcher}
\date{March 2026}

\begin{document}

\maketitle
\begin{center}
\large Paper 25.3
\end{center}

\begin{abstract}
This paper freezes the full HAOS-IIP no-distance combinatorial austerity line as a completed C-line document. It carries the line from Stage C0.1 through Stage C0.20, includes the Stage 11 integrated austerity-resolution probe, and closes the branch with Stage 12, a label-free structural closure test. The line begins as a strict no-distance control: coordinate distance is removed from the interaction kernel and replaced by graph-internal structure on an explicit edge-graph control branch. It then expands in four directions. The first direction, Stages C0.5--C0.9, asks whether combinatorial survivor rules, incidence-seeded mismatch accumulation, and softer bounded-path kernels can produce sequencing, explicit kernel inconsistency, or broad irreversibility. The second direction, Stages C0.10--C0.16, adds a static harmonic-address compatibility factor on the clustered Dirac--Kahler seed and asks whether topology selection sharpens into a resonance ladder, a detuning phase boundary, hysteretic return error, selector-withdrawal retention, or a locking mechanism. The third direction tests whether a signed combinatorial Dirac--Kahler candidate can be constructed directly from C0 primitives and whether any part of the harmonic detuning and protection structure survives on that derived complex. The fourth direction starts from that bridge scaffold and asks the fairest late questions: does a recovered bridge braid window survive selector withdrawal, does it leave an armed-state hysteresis witness, do the open bridge obligations converge when mismatch orientation, signed evolution, harmonic-address sectoring, and topology survival are measured synchronously, and does any of that survive a label-free bidirectional closure test. The final result is bounded and negative in the exact way a control program should be. Graph-structural class survival survives metric removal, incidence-forced protected sequencing appears reproducibly under external combinatorial forcing, the harmonic-address branch shows a weak but real topology-selection signal with a sharp detuning threshold, the bridge validation closes the signed chain identities exactly on the tested representative graph complexes, and corrected bridge hysteresis probes open clustered armed-state irreversibility witnesses. But selector withdrawal still fails to retain topology, Stage 11 splits rather than resolves the open bridge obligations, and Stage 12 closes negatively: label-free cycling preserves a tight evolution surrogate and only a narrow irreversibility signal, while orientation fails, intrinsic grading fails, and family-wide recovery fails. The final boundary is therefore explicit. The full C line supports graph-structural class survival, limited protected sequencing, weak harmonic-address topology selection, clustered armed-state irreversibility on the bridge scaffold, and a real C0-to-DK bridge candidate, but it does not support topology locking, intrinsic grading, or family-wide recovery on the no-distance branch.
\end{abstract}

\section{Introduction}

The C0 line is a bounded control program inside HAOS-IIP, not a replacement core. Its purpose is narrower:
\begin{quote}
\emph{What survives once coordinate distance is removed, and what further combinatorial ingredients are still insufficient to turn selection into irreversibility or locking?}
\end{quote}

The early answer was reported in Paper 24.1, which froze Stages C0.1 through C0.4 as a metric-removal and topology-control sequence. Paper 25.1 then carried the line through C0.20 and Stage 11. The present paper is the closure version of that arc. It incorporates Stage 12 as the decisive label-free austerity test and freezes the combinatorial branch as a completed C line while keeping the core/model boundary explicit:
\begin{itemize}
\item the frozen core remains the HAOS-IIP axioms of recoverability, interaction primacy, representational regimes, viability-constrained dynamics, and invariant residue;
\item the C0 line remains a regime-specific control branch beneath that core;
\item positive and negative numerical outcomes are treated as branch boundaries, not as revisions of the axioms.
\end{itemize}

\section{Core Boundary and Spectral Framing}

The recent canonical cleanup matters here. On Laplacian-style branches, HAOS recoverability can be written as an operator-level spectral condition rather than as free-floating language. If
\[
L = D-A, \qquad 0=\lambda_0 \le \lambda_1 \le \lambda_2 \le \cdots,
\]
then the cleanest scalar-branch compression is
\[
e^{-\lambda_1 T}\delta < \epsilon_c,
\]
where $\delta$ is perturbation size, $T$ is the admissible recovery horizon, and $\epsilon_c$ is the viable coherence threshold. In words: the low end of the interaction spectrum must damp perturbations back below the viable coherence band quickly enough to restore operational structure.

This sharpens the language of recoverability on Laplacian branches, but it does not rewrite any C0 outcome. It clarifies how the scalar control line should be read:
\begin{itemize}
\item C0 positives remain positives only if recoverability improves inside the branch;
\item C0 nulls remain nulls even if a branch looks visually structured;
\item a partial opening of kernel inconsistency is not the same thing as broad irreversibility;
\item topology selection is not the same thing as topology locking.
\end{itemize}

\section{Design Discipline Across the C0 Line}

The C0 line actually contains two sub-branches that share a common discipline but answer different questions.

\subsection{C0.1--C0.9: pure no-distance control branch}

These stages use an explicit projected-transverse edge-graph control operator. The graph size, representative encounter families, and shell-weight measurement vocabulary are fixed while the kernel or survivor rule is varied using only graph-internal quantities:
\begin{itemize}
\item adjacency,
\item bounded path length in edge count,
\item local degree,
\item local motif incidence,
\item combinatorial Laplacian structure.
\end{itemize}

The three representative families are:
\begin{itemize}
\item \texttt{clustered\_composite\_anchor},
\item \texttt{counter\_propagating\_corridor},
\item \texttt{phase\_ordered\_symmetric\_triad}.
\end{itemize}

\subsection{C0.10--C0.14: harmonic-address extension}

These stages shift to the clustered Phase III Dirac--Kahler seed while keeping the projected-transverse sector, graph resolution, and topology measurement stack fixed. A static harmonic-address compatibility factor is added as a relational multiplier on the combinatorial operator. The branch therefore asks a different question:
\begin{quote}
\emph{Can harmonic-address compatibility select, stratify, or lock topology without metric distance?}
\end{quote}

\section{Branch A: Metric Removal and Topology Controls (C0.1--C0.4)}

Paper 24.1 already froze the first four stages. The present synthesis keeps that result intact:
\begin{itemize}
\item C0.1: primary persistence-class behavior survives removal of coordinate-distance weighting on the explicit edge-graph control branch, although transport texture shifts.
\item C0.2: degree-spectrum variation reshapes transport span and flow concentration but does not overturn the coarse topology labels.
\item C0.3: the no-distance branch does not produce a clean braid baseline; the nearest honest surrogate, the counter-propagating corridor, remains \texttt{smeared\_transfer} under all tested incidence-noise perturbations.
\item C0.4: small local motifs do not seed a reproducible new topology class; eight of nine runs are neutral and the only non-neutral response is an isolated diamond-induced channel-splitting case.
\end{itemize}

The stable conclusion from C0.1--C0.4 therefore remains:
\begin{quote}
\emph{The no-distance branch supports graph-structural class survival and transport reshaping, but not a reproducible combinatorial braid family or a motif-seeded topology mechanism.}
\end{quote}

\section{Branch B: Sequencing, Mismatch, and Kernel Inconsistency (C0.5--C0.9)}

\subsection{C0.5: static sequencing baseline}

Stage C0.5 replaces the assumed packet-evolution clock with a combinatorial sequencing rule built from persistence-class stabilization, incidence-mismatch accounting, kernel-update consistency, and protected topology-stability blocks. The result is mostly static. Eight of nine runs freeze into \texttt{static\_regime}, the mismatch ledger stays zero everywhere, and only one protected ordering event appears on the symmetric triad.

This is a useful negative baseline: the unforced branch is too rigid to generate sequencing internally.

\subsection{C0.6: incidence-seeded mismatch accumulation}

Stage C0.6 deliberately injects the noise and motif menus from C0.3 and C0.4 after the protected stability constraint is activated. Here the branch opens clearly:
\begin{itemize}
\item positive mismatch density appears in $9/9$ runs,
\item protected arrow forcing appears in $6/9$ runs,
\item the protected family labels remain intact: three \texttt{trapped\_local}, three \texttt{smeared\_transfer}, three \texttt{split\_channel\_exchange}.
\end{itemize}

This is the strongest positive result in the sequencing line. The branch does support incidence-forced protected ordering. But the mechanism is still not inconsistency-mediated because kernel-consistency failures remain zero.

\subsection{C0.7: original alternation remains too rigid}

Stage C0.7 adds explicit adjacency-versus-graph-shell alternation on top of the positive-mismatch runs from C0.6. The result is cleanly negative:
\begin{itemize}
\item explicit consistency failures remain $0/9$,
\item protected arrow under inconsistency remains $0/9$,
\item the general protected ordering inherited from C0.6 survives only as an ordering-without-inconsistency effect.
\end{itemize}

So the original survivor-kernel pair is too rigid to support explicit inconsistency as the driver.

\subsection{C0.8 and C0.9: softer survivor maps partially open, then fail to scale}

Stage C0.8 replaces the rigid graph-shell survivor map with a bounded-path kernel weighted by $1/(h+1)$ up to a fixed hop cap and upgrades motifs to bi-root and two-layer injections. This is the first stage that opens explicit inconsistency:
\begin{itemize}
\item explicit consistency failures appear in $5/9$ runs,
\item protected-arrow events under weaker inconsistency appear in $2/9$ runs,
\item the hybrid noise-plus-motif protocol is the only forcing family that produces protected ordering under explicit inconsistency, with positive cases on the clustered anchor and the symmetric triad.
\end{itemize}

Stage C0.9 tries to scale that partial opening with tri-root and three-layer motifs. It fails:
\begin{itemize}
\item explicit consistency failures drop to $4/9$,
\item tri-root-specific failures appear in only $2/9$,
\item protected-arrow events under weaker or tri-root inconsistency fall to $0/9$,
\item the irreversibility scaling factor is $0.0$ relative to the C0.8 partial positive.
\end{itemize}

\begin{table}[t]
\centering
\scriptsize
\begin{tabular}{@{}lcccc@{}}
\toprule
Stage & Positive mismatch & Explicit inconsistency & Protected arrow & Main read \\
\midrule
C0.5 & $0/9$ & $0/9$ & $1/9$ & mostly static sequencing baseline \\
C0.6 & $9/9$ & $0/9$ & $6/9$ & incidence-forced protected ordering \\
C0.7 & $9/9$ & $0/9$ & $0/9$ & original alternation too rigid \\
C0.8 & $9/9$ & $5/9$ & $2/9$ & partial inconsistency opening under weaker kernel \\
C0.9 & $9/9$ & $4/9$ & $0/9$ & tri-root scaling fails \\
\bottomrule
\end{tabular}
\caption{Sequencing and inconsistency block. The branch supports reproducible mismatch accumulation and a partial inconsistency opening under a softer bounded-path survivor map, but not broad irreversibility scaling.}
\label{tab:c0b}
\end{table}

The correct combined boundary for C0.5--C0.9 is therefore:
\begin{quote}
\emph{The no-distance branch supports incidence-forced protected sequencing and can be pushed into explicit inconsistency under softer survivor maps, but it does not support a broad claim that irreversibility follows automatically from that inconsistency.}
\end{quote}

\section{Branch C: Harmonic-Address Selection, Detuning, Hysteresis, and Selector Withdrawal (C0.10--C0.16)}

\subsection{C0.10: weak proto-selection, no persistence gain}

Stage C0.10 adds a static harmonic-address compatibility factor to the clustered Dirac--Kahler seed. This opens the first clean topology-selection signal:
\begin{itemize}
\item exact-match runs preserve the braid class in $2/9$ runs,
\item near-match and incompatible sectors relax to \texttt{transfer\_smeared},
\item persistence gain remains exactly $0.0$ across the matrix.
\end{itemize}

So harmonic addressing acts as a weak proto-selection rule, not as a persistence generator.

\subsection{C0.11 and C0.12: ladder plus threshold}

Stage C0.11 asks whether discrete harmonic-address differences generate a resonance ladder. The answer is yes, but in a split sense:
\begin{itemize}
\item topology shows two visible bands: same-address versus all nonzero address differences,
\item resonance alignment further stratifies the nonzero sector into $d1/d4$ and $d2/d3$.
\end{itemize}
The same-address braid sector is not fully refinement-stable, so the ladder is stronger in alignment than in topology stability.

Stage C0.12 then scans a denser detuning ladder. The result is a sharp threshold:
\begin{itemize}
\item \texttt{braid\_like\_exchange} persists through $\delta = 0.625$ in $6/9$ runs,
\item the topology flips once near $\delta = 0.750$,
\item the remaining three runs are \texttt{transfer\_smeared} for $\delta = 0.750, 0.875, 1.000$.
\end{itemize}

\subsection{C0.13 and C0.14: boundary refinement and locking null}

Stage C0.13 resolves the detuning corridor near the observed transition. The local picture is sharp:
\begin{itemize}
\item in the baseline, motif-weak-minus, and beta-weak families, \texttt{address\_protected\_braid} survives at $\delta = 0.66, 0.69$ and flips to \texttt{transfer\_smeared} at $\delta = 0.72$,
\item motif-weak-plus, degree-mild, and degree-asymmetric collapse the protected window before the scan even enters the corridor.
\end{itemize}

The global verdict remains \texttt{fragmented multi-window protection} because no family satisfies the strict three-consecutive-sample resolution rule. The boundary is visually sharp in the baseline-like families but not formally resolved by the atlas criterion.

Stage C0.14 then tests the next honest escalation: can bounded local latches convert topology selection into locking? The answer is no.
\begin{itemize}
\item all $9/9$ runs remain \texttt{no\_locking\_effect},
\item no \texttt{locked\_braid} or \texttt{quasi\_bounded\_exchange} class appears,
\item the exact-match family stays \texttt{braid\_like\_exchange},
\item the boundary and fragile families stay \texttt{transfer\_smeared},
\item the harmonic phase latch actually shortens braid survival on the exact-match seed by $0.1253$.
\end{itemize}

\begin{table}[t]
\centering
\scriptsize
\begin{tabular}{@{}lccc@{}}
\toprule
Stage & Main class counts & Strongest positive & Boundary \\
\midrule
C0.10 & $2$ braid, $7$ smear & exact-match address preserves braid & no persistence gain \\
C0.11 & $2$ braid, $7$ dispersive & resonance ladder in alignment & same-address braid not refinement-stable \\
C0.12 & $6$ braid, $3$ smear & sharp flip near $\delta = 0.750$ & not a smooth decay and not multi-plateau \\
C0.13 & $6$ protected braid, $36$ smear & narrow local protected window & no family satisfies strict 3-sample rule \\
C0.14 & $3$ braid, $6$ smear & none & all $9/9$ runs are \texttt{no\_locking\_effect} \\
\bottomrule
\end{tabular}
\caption{Harmonic-address block. Harmonic addressing selects and stratifies topology, but the branch still closes on a locking null result.}
\label{tab:c0c}
\end{table}

The correct combined boundary for C0.10--C0.14 is:
\begin{quote}
\emph{Harmonic addressing produces weak topology selection and a real detuning threshold on the clustered seed, but the tested bounded latch rules do not convert that selection into locking.}
\end{quote}

\subsection{C0.15 and C0.16: hysteresis and selector-withdrawal do not rescue locking}

The next two stages ask the sharper derivational questions that C0.14 leaves open.

Stage C0.15 replaces one-way detuning scans with a forward/reverse hysteresis protocol. The same harmonic-address family is driven out and back through the detuning corridor while the return error and hysteresis area are measured directly. The result is a partial positive only:
\begin{itemize}
\item $4/9$ runs classify as \texttt{irreversible\_witness},
\item $5/9$ runs classify as \texttt{hysteretic\_retrace},
\item the strongest non-return cases are the exact plain, exact motif, fragile degree, and fragile motif protocols,
\item but broad irreversibility is not established.
\end{itemize}
So the branch does support path dependence, but not a strong claim that irreversibility now follows generically from harmonic-address forcing.

Stage C0.16 then asks the cleaner locking question: once a topology has been selected, does it remain after the selector is removed? The answer is a clean no.
\begin{itemize}
\item $0/9$ runs produce any retained locking class,
\item $3/9$ runs classify as \texttt{selector\_dependent\_collapse},
\item $6/9$ runs classify as \texttt{no\_locking\_after\_withdrawal},
\item the threshold-protected family can still arm into braid under selection, but immediately collapses when the selector is withdrawn.
\end{itemize}

These two stages sharpen the earlier locking null rather than overturning it. The honest extension of the harmonic-address boundary is therefore:
\begin{quote}
\emph{The harmonic-address branch supports weak topology selection, a real detuning threshold, and partial hysteresis, but selected topology does not survive selector withdrawal and therefore still does not qualify as locked.}
\end{quote}

\section{C0-to-DK Bridge Validation and Late Bridge Extensions}

The bridge note is not a new persistence claim. It is a validation scaffold that asks whether a signed combinatorial Dirac--Kahler candidate can be built directly from C0 primitives and whether any part of the late C0 selection/protection structure survives on that derived complex.

The first important result is algebraic rather than phenomenological. On the three representative graph families
\begin{itemize}
\item \texttt{clustered\_composite\_anchor},
\item \texttt{counter\_propagating\_corridor},
\item \texttt{phase\_ordered\_symmetric\_triad},
\end{itemize}
the signed bridge complex closes exactly:
\[
\|d_1 d_0\| = 0, \qquad \|D^2 - \Delta\| = 0, \qquad \|d_0^{T} d_0 - L_0\| = 0.
\]
So the bridge candidate is not just a verbal analogy. It is a real signed chain complex on the tested graph families.

The second result is selective and narrower. On the derived complex:
\begin{itemize}
\item the clustered family opens a measurable braid-to-smear bridge corridor, with \texttt{braid\_like\_exchange} at $\delta = 0.0, 0.125$ and \texttt{transfer\_smeared} from $\delta = 0.25$ onward;
\item the focused clustered protection panel reproduces the full tested qualitative pattern: balanced baseline braid, phase-edge smear, width-asymmetry smear, degree-skew tolerance, and zero-form fragility;
\item the corridor and triad families do not recover a broad braid sector.
\end{itemize}

The bridge note therefore earns a bounded but real statement:
\begin{quote}
\emph{A combinatorial Dirac--Kahler candidate can be constructed directly from C0 primitives, and its signed chain identities close exactly on the tested representative graph complexes. But orientation source, evolution-law equivalence, intrinsic grading source, and family-wide protection recovery remain conditional.}
\end{quote}

This is useful because it tightens the HAOS-IIP core/model boundary rather than weakening it. The graded sector is no longer only an imported Phase III object; there is now a real C0-derived bridge scaffold. But it remains a scaffold, not a finished derivation.

\subsection{C0.17--C0.20: fair withdrawal stays negative, armed-state hysteresis opens irreversibility}

The first two bridge-local latch stages, C0.17 and C0.18, were useful but not yet fair enough as retention tests. Both were fully negative, but they mixed failure-to-arm with failure-to-retain: the clustered bridge family often left the braid class during arming itself.

Stage C0.19 corrects that by withdrawing from the real recovered clustered bridge braid window rather than from a fixed arming cutoff. The corrected read is stronger, not weaker:
\begin{itemize}
\item the clustered bridge family does recover a real braid window at steps $138$--$142$,
\item all three fair withdrawal protocols then collapse that armed state to \texttt{unresolved\_mixed},
\item post-withdrawal braid dwell is only $1.0$ step in every clustered run,
\item the best label is \texttt{topology\_afterglow}, not locking.
\end{itemize}
So the late locking null is no longer vulnerable to the objection that the bridge stages withdrew from the wrong part of the trajectory. Even when the selector is removed from a real bridge braid window, topology does not hold itself together.

Stage C0.20 then asks the irreversibility question from that same armed state. This is the strongest late positive now available on the bridge scaffold:
\begin{itemize}
\item all three clustered armed-state forward/reverse protocols finish in \texttt{transfer\_smeared} even after exact-match restoration,
\item all three clustered runs therefore classify as \texttt{irreversible\_witness},
\item the strongest protocol is the frozen-hold cycle, with hysteresis areas $0.1341$ in flow and $0.0977$ in loop score,
\item corridor and triad still fail to open armed braid windows on the bridge scaffold and therefore remain \texttt{no\_armed\_braid\_window}.
\end{itemize}

The combined bridge extension is therefore asymmetric:
\begin{quote}
\emph{The bridge scaffold now supports explicit armed-state irreversibility on the clustered family, but even a fair selector-withdrawal test does not produce topology locking.}
\end{quote}

\section{Stage 11 Integrated Austerity Resolution}

Stage 11 is the first convergence probe rather than a narrow micro-control. It forces four previously open structures to co-evolve synchronously on matched graph families:
\begin{itemize}
\item mismatch-ledger orientation,
\item signed combinatorial field evolution,
\item intrinsic harmonic-address sectoring,
\item topology survival.
\end{itemize}
The point is to test whether any positive signal still survives once these structures are measured in the same state history rather than across separate runs.

The result is globally negative, but in a sharpened way:
\begin{itemize}
\item \texttt{O\_C0\_ORIENT} is positive: two families show persistent positive ledger alignment, with medians $0.9996$ on the clustered seed and $0.9937$ on the triad;
\item \texttt{O\_C0\_EVOLUTION} is positive: the coarse/fine evolution mismatch ratio remains tightly bounded around $1.0056$;
\item \texttt{O\_C0\_GRADING} is negative: sector entropy drops only weakly, from $0.2114$ to $0.2091$, and does not resolve into a convincing intrinsic-grade structure;
\item \texttt{O\_C0\_RECOVERY} is negative: the required ordering \texttt{clustered} $\ge$ \texttt{corridor} $\ge$ \texttt{triad} fails in every setting, so the ordered-setting fraction is $0.0000$.
\end{itemize}

So Stage 11 does not collapse the bridge obligations into a single structural win. It splits them:
\begin{quote}
\emph{Mismatch-orientation and signed-field evolution now look genuinely coupled, but harmonic address still behaves more like selector phenomenology than intrinsic grading, and family-wide recovery remains unresolved.}
\end{quote}

\section{Stage 12 Intrinsic Grading and Structural Irreversibility Closure}

Stage 12 is the closure test rather than another phenomenology scan. It removes explicit harmonic-address labels from the evolution rule, forces forward/reverse/restoration cycles on the same common lattice, applies the same protocol to the clustered, corridor, and triad families, and doubles the resolution to ask the hardest late-C-line question:
\begin{quote}
\emph{Does the combinatorial line close as an intrinsic structural framework once selector dressing is removed, or does it remain a bounded selector-rich bridge?}
\end{quote}

The core run lattice is the smallest one that still closes that logic:
\begin{itemize}
\item three families: clustered protected seed, corridor weak-protection seed, triad competition seed;
\item three protocols: forward only, forward/reverse/restore, forward/reverse/re-perturb;
\item two resolutions: base and refinement doubling.
\end{itemize}
This gives eighteen core runs under a single label-free rule set.

The result is \texttt{negative\_closure}. Two strengthened obligations survive; two do not.
\begin{itemize}
\item \texttt{O\_C0\_EVOLUTION} survives the harsher setting: the median refinement ratio remains $1.0381$, so mismatch-step projection and signed-field flux stay tightly aligned under label-free cycling.
\item Structural irreversibility survives only in a narrow sense: the bidirectional positive fraction is $0.3333$.
\item \texttt{O\_C0\_ORIENT} fails once explicit addressing is removed: the positive-family count falls to $0$.
\item Intrinsic grading fails: entropy remains $0.2570$, separation stays modest at $0.2019$, and detuning monotonicity turns negative at $-0.3616$.
\item Family-wide recovery fails completely: the ordered-setting fraction is $0.0000$.
\end{itemize}

\begin{table}[t]
\centering
\scriptsize
\begin{tabular}{@{}lcc@{}}
\toprule
Diagnostic & Value & Read \\
\midrule
\texttt{O\_C0\_ORIENT} positive-family count & $0$ & fails under label-free cycling \\
\texttt{O\_C0\_EVOLUTION} median refinement ratio & $1.0381$ & survives \\
Recovered grading entropy & $0.2570$ & too high for closure \\
Recovered grading separation & $0.2019$ & weak structuralization \\
Detuning monotonicity & $-0.3616$ & wrong sign for intrinsic grading \\
Family recovery ordered fraction & $0.0000$ & fails \\
Bidirectional positive fraction & $0.3333$ & limited irreversibility only \\
\bottomrule
\end{tabular}
\caption{Stage 12 closure diagnostics. Label-free cycling preserves the evolution surrogate and a narrow irreversibility signal, but not orientation, intrinsic grading, or family-wide recovery.}
\label{tab:stage12}
\end{table}

The family structure is also informative. The strongest irreversibility signal shifts into the corridor family rather than staying on the clustered bridge scaffold:
\begin{itemize}
\item clustered medians: alignment $0.1161$, protected fraction $1.0000$, irreversibility $0.0426$, entropy $0.2458$;
\item corridor medians: alignment $0.0339$, protected fraction $0.5541$, irreversibility $1.0503$, entropy $0.3501$;
\item triad medians: alignment $0.1081$, protected fraction $0.9965$, irreversibility $0.0269$, entropy $0.2305$.
\end{itemize}
So Stage 12 does not reward the most protected family. It exposes that the surviving irreversibility signal is narrow, protocol-sensitive, and not the same thing as a universal protection hierarchy.

The clean closure statement is therefore:
\begin{quote}
\emph{Label-free bidirectional cycling preserves a tight evolution surrogate and a narrow irreversibility signal, but it does not recover intrinsic grading, does not stabilize the orientation ledger, and does not produce a universal protection ordering across families.}
\end{quote}

\section{What Changes and What Does Not}

The updated paper changes the map of the C0 line, but not the core or the main phase boundaries.

\subsection{What changes}

\begin{itemize}
\item The C0 line is no longer just a metric-removal control. It now includes a full sequencing/inconsistency block, a harmonic-address selection/locking block, a late bridge irreversibility block, and a synchronous integrated-resolution probe.
\item The strongest late positive is no longer ``metric removal preserves some classes.'' It is C0.6: incidence-forced protected sequencing with positive mismatch in $9/9$ runs and protected arrow forcing in $6/9$ runs.
\item The strongest late negative is no longer just ``motifs do not seed braid.'' It is now the combined C0.14--C0.19 boundary: selection does not become locking under bounded latches, selected topology does not survive selector withdrawal, and even fair bridge withdrawal from a real braid window yields only topology afterglow.
\item A new bridge-level positive now exists at two levels: the C0-to-DK signed complex closes exactly on the tested representative graph families and recovers a clustered-only detuning/protection signal, and C0.20 later adds an explicit armed-state irreversibility witness on that same clustered bridge scaffold.
\item Stage 11 and Stage 12 now split the obligation map cleanly. Orientation and evolution couple enough to survive the integrated Stage 11 test, but Stage 12 shows that only evolution survives the harsher label-free closure probe. Intrinsic grading and family-wide recovery remain negative, and the surviving irreversibility signal is narrow rather than universal.
\item The combinatorial line is now actually closed. Stage 12 does not open a new rescue path; it freezes the branch as a bounded negative/partial result with an explicit final boundary.
\end{itemize}

\subsection{What does not change}

\begin{itemize}
\item The frozen HAOS-IIP core axioms do not change.
\item Paper 24.1 remains the frozen early synthesis for C0.1--C0.4.
\item Paper 25.1 remains the Stage-11 synthesis release; the present paper is the closure release rather than a replacement of its internal results.
\item Phase I remains closed by the same boundary: frozen pre-geometry can organize, but it does not bind.
\item Phase II is still not rescued by the C0 line. C0.8 opens only a partial inconsistency signal and C0.9 fails to scale it.
\item Phase III topology findings remain topology findings; the harmonic-address branch adds a selector on that clustered seed but not a persistence law, and neither the bridge extensions nor Stage 11 collapse the gap between C0 and the full Phase III graded sector.
\end{itemize}

\section{Conclusion}

The full C0.1--C0.20 line, together with Stage 11 and the decisive Stage 12 closure probe, now has a frozen architecture and a final bounded read.

First, metric removal does not erase the no-distance branch. Real graph-structural class survival remains after coordinate distance is removed. Degree landscape, incidence noise, and motif forcing reshape transport texture and local class behavior, but they do not magically produce a robust braid family.

Second, sequencing and inconsistency are not automatic. The branch supports reproducible mismatch accumulation and protected ordering under external forcing, and a weaker bounded-path survivor map can open explicit kernel inconsistency. But that opening is partial and does not scale into a broad irreversibility claim. Later hysteresis probes show path dependence in part of the harmonic-address matrix, bridge-local hysteresis later opens a real clustered armed-state irreversibility witness, and Stage 12 still shows that the surviving bidirectional irreversibility signal is narrow rather than universal.

Third, harmonic addressing behaves like a selector, not a locker. Exact-match sectors preserve braid on the clustered seed, resonance alignment stratifies into discrete address bands, and a narrow detuning threshold appears near $\delta \approx 0.75$ on the original clustered DK branch. But bounded latches do not convert that selection into a locked or quasi-bounded exchange regime, selector withdrawal collapses the selected topology instead of retaining it, and even fair bridge withdrawal from a real recovered braid window yields only one-step afterglow.

Fourth, the graded bridge is now real but still conditional. A C0-derived signed Dirac--Kahler candidate closes exactly on the tested representative graph complexes and recovers a clustered-only detuning/protection signal. From that scaffold, armed-state bridge hysteresis yields an explicit clustered irreversibility witness. But Stage 11 and Stage 12 together show that intrinsic grading and family-wide recovery do not converge, and Stage 12 further shows that the orientation ledger does not survive the full label-free closure test. That is enough to justify the bridge as a scaffold with a real late irreversibility signal, not enough to claim that the graded sector is fully derived from C0 alone.

The final C0 boundary can therefore be stated without hype:
\begin{quote}
\emph{The no-distance combinatorial line supports graph-structural class survival, limited protected sequencing, weak harmonic-address topology selection, explicit clustered armed-state irreversibility on the bridge scaffold, and a real C0-to-DK bridge candidate, but it does not support topology locking, intrinsic grading, or family-wide recovery on the no-distance branch.}
\end{quote}

That is not a failure of the line. It is exactly the kind of bounded map a control program is supposed to produce, and it is why the combinatorial austerity line can now be treated as closed.

\end{document}
